\documentclass[11pt, a4paper]{ctexart}
\usepackage{geometry}
\geometry{left=2cm, right=2cm, top=2.5cm, bottom=2.5cm}
\usepackage{hyperref}
\usepackage{color}
\usepackage{xcolor}
\usepackage{booktabs}
\usepackage{enumitem}
\usepackage{tcolorbox}
\usepackage{longtable}
\usepackage{amsmath}
\usepackage{amssymb}
\usepackage{graphicx}
\usepackage{fancyhdr}

\hypersetup{
    colorlinks=true,
    linkcolor=blue,
    filecolor=magenta,      
    urlcolor=cyan,
}

% Colors
\definecolor{primary}{HTML}{1A365D}
\definecolor{secondary}{HTML}{2B6CB0}
\definecolor{lightbg}{HTML}{F7FAFC}
\definecolor{border}{HTML}{E2E8F0}

% Header/Footer
\pagestyle{fancy}
\fancyhf{}
\fancyhead[L]{\textcolor{gray}{PowerRAG 期末汇报精简演讲手稿 (10分钟陈述专用)}}
\fancyhead[R]{\textcolor{gray}{\leftmark}}
\fancyfoot[C]{\thepage}
\renewcommand{\headrulewidth}{0.4pt}
\renewcommand{\footrulewidth}{0pt}

% Section formatting
\ctexset{
  section = {
    format = \Large\bfseries\color{primary},
    aftername = \hspace{0.5em}
  },
  subsection = {
    format = \large\bfseries\color{secondary},
    aftername = \hspace{0.5em}
  }
}

\begin{document}

\section*{3 分钟压缩版}
\addcontentsline{toc}{section}{3 分钟压缩版}

各位老师好，我汇报的项目叫 PowerRAG。\\

这个项目的起点不是做一个普通问答网页，而是来自我们之前鸿雁计划里的研究设想：面向动力装备故障诊断，构建知识图谱增强的 RAG 框架、评测指标和中文 benchmark。课程阶段我做的 PowerRAG，可以理解为这个设想的第一版工程化落地。\\

我们面对的不是一道问答题，而是一批真实资料。动力装备资料里有 PDF、扫描件、JSON 标注、维修报告、页码、表格和图像信息。如果直接交给大模型，它可能能说出一段答案，但很难说明答案来自哪里，也很难被检查。\\

所以项目第一步是资料入口。JSON 走结构化解析，文字 PDF 走文本抽取，扫描 PDF 先走 OCR。这样不同资料才能进入同一套后续流程。\\

第二步是证据构建。我们不是简单把文本塞进向量库，而是把资料切成带来源、页码、标题和上下文的证据块。这样系统回答时，答案可以回到原始资料，而不是只生成一句看起来合理的话。\\

第三步是检索实验。我们用 60 道评测题比较了 keyword、hybrid\_rrf 和 dense\_hashing。结果 keyword 的 question recall@K 是 0.833，hybrid\_rrf 是 0.817，dense\_hashing 是 0.583。这说明在动力装备资料中，专业术语和型号本身就是强检索信号，复杂方法不一定天然更好，必须用评测决定路线。\\

第四步是 GraphRAG。普通 RAG 适合回答有直接证据的问题，但故障、部件、工况、参数之间存在关系链，所以我们引入 GraphRAG，把实体和关系组织起来，用于关系追踪、社区摘要和全局综合。但它不是替代所有 RAG，而是根据问题类型进行增强。\\

最后是评测闭环。系统不能只会演示，还要能评分、能比较 baseline、能分析失败原因。这里要分清两类结果：动力装备 60 题 baseline 用来判断本项目检索路线；公开 benchmark 外部评测则用来对标成熟 RAG 任务。我们下载并接入了 legal-rag-bench 和 ragbench-hotpotqa-validation，各取 50 个 case，最终 completion audit 总分达到 93.47 / 100。失败可能来自 OCR、分块、召回、排序、图谱关系或证据绑定。评测的意义，就是把“效果不好”变成下一轮可以执行的优化任务。\\

总结来说，PowerRAG 是一个面向动力装备资料的可信问答与评测平台。它从资料进入开始，到证据构建、检索路由、GraphRAG 增强，再到评测归因，形成了一条可以继续迭代的工程链路。下一阶段会继续完善 embedding、reranker、GraphRAG 同题评测和更大规模中文 benchmark。\\


\noindent\rule{\textwidth}{0.5pt}


\section*{10 句背诵提纲}
\addcontentsline{toc}{section}{10 句背诵提纲}

\begin{enumerate}[leftmargin=*]
\item PowerRAG 不是普通问答网页，而是面向动力装备资料的可信知识系统。
\item 它延续了鸿雁计划里“知识图谱增强工业故障诊断 RAG”的立项初心。
\item 我们面对的是 PDF、扫描件、JSON、维修报告和标准文档，不是一道干净问答题。
\item 第一关是资料入口：JSON 解析，PDF 抽取，扫描件 OCR。
\item 第二关是证据构建：chunk 必须保留来源、页码、标题和上下文。
\item 检索实验说明，动力装备资料里的专业术语让 keyword 在当前 60 题中表现很强。
\item 所以方法不是越复杂越好，而是要用评测决定路线。
\item GraphRAG 的价值在关系追踪、社区摘要和全局综合，不是替代所有 RAG。
\item 评测平台让项目从“能演示”变成“能评分、能归因、能优化”，公开 benchmark 外部评测达到 93.47 / 100。
\item 下一步是更正式的 embedding、reranker、GraphRAG 同题评测和更大规模中文 benchmark。
\end{enumerate}

\noindent\rule{\textwidth}{0.5pt}


\section*{答辩追问备用回答}
\addcontentsline{toc}{section}{答辩追问备用回答}

\subsection*{如果老师问：你们和普通 RAG 有什么区别？}
\addcontentsline{toc}{subsection}{如果老师问：你们和普通 RAG 有什么区别？}

普通 RAG 更像是把文档切片后做检索问答。我们这里更强调真实资料工程和可信闭环：资料入口要处理 JSON、PDF、扫描件和 OCR；证据块要保留来源和页码；检索方法要用评测比较；GraphRAG 用来补关系型和全局综合问题；最后还要有评测平台做 baseline 和失败归因。所以区别不在于有没有调用大模型，而在于有没有把资料、证据、图谱和评测组织成一个系统。\\

\subsection*{如果老师问：为什么 keyword 比 dense 还高？}
\addcontentsline{toc}{subsection}{如果老师问：为什么 keyword 比 dense 还高？}

这个结果和领域资料特点有关。动力装备资料里有很多强术语，比如设备名、部件名、故障名、型号和工况参数。这类词本身就是非常明确的检索锚点。当前 dense\_hashing 可能没有充分理解这些专业术语，所以表现不如 keyword。这不代表语义检索没有价值，而是说明后续需要更好的 embedding、reranker 和问题路由。\\

\subsection*{如果老师问：GraphRAG 到底解决什么问题？}
\addcontentsline{toc}{subsection}{如果老师问：GraphRAG 到底解决什么问题？}

GraphRAG 主要解决普通文本片段难以表达的关系型问题。比如故障和部件、工况和参数、维修动作和原因之间的连接。它适合关系追踪、社区摘要和跨文档综合。但对于简单定义或事实问题，普通 RAG 可能更直接，所以我们把 GraphRAG 定位为增强路线，不是替代路线。\\

\subsection*{如果老师问：评测平台现在成熟吗？}
\addcontentsline{toc}{subsection}{如果老师问：评测平台现在成熟吗？}

当前是第一版系统化评测闭环，已经可以围绕动力装备 60 题做 baseline 比较和失败分析；同时也能接入下载到本地的公开 RAG benchmark。外部评测里，\texttt{open\_source\_90\_completion\_audit\_50x2\_auto} 这次运行覆盖 100 个 case，总分 93.47 / 100，其中 legal-rag-bench 为 96.68，ragbench-hotpotqa-validation 为 90.25。这里需要说明的是，这代表 retrieval-first 外部 benchmark 质量门禁通过，不等于最终动力装备专家 benchmark 已经完成。按照鸿雁计划，后续目标是扩展到 100 到 200 道以上的中文动力装备故障诊断评测题，并引入更完整的指标和专家评分。\\

\subsection*{如果老师问：93.47 分能说明什么？}
\addcontentsline{toc}{subsection}{如果老师问：93.47 分能说明什么？}

它说明评测系统已经开始发挥工程作用。我们不是凭主观感觉说系统变好了，而是把公开 benchmark 下载到本地，用同一套评测入口反复检查解析、切分、检索、召回和重排策略。最后在 100 个公开 benchmark case 上取得 93.47 / 100，并通过 Open Source 90 External Benchmark Gate。但我不会把它说成最终工业验收结果，因为动力装备真实语料、专家 gold set 和更大规模中文 benchmark 还需要继续补充。\\

\subsection*{如果老师问：这个项目下一步最重要的是什么？}
\addcontentsline{toc}{subsection}{如果老师问：这个项目下一步最重要的是什么？}

我认为最重要的是把评测闭环继续做实。具体包括更换正式 embedding，引入 reranker，对 GraphRAG local 和 global 做同题对比，扩展评测集，并把失败归因和任务队列打通。这样项目才能从课程演示继续走向可复现实验平台。\\


\end{document}
